% Introduction — Agentic EEG preprocessing pipeline with FOOOF-based signal quality optimization
% Target journal: Journal of Neural Engineering (methods-oriented, neuroengineering audience)
% All \cite{} keys validated against step4_references.bib (53 entries)

\section{Introduction}
\label{sec:intro}

\subsection{Context and motivation}

Electroencephalography (EEG) is among the most widely used non-invasive techniques for studying human brain function, but raw EEG is heavily contaminated by ocular, muscular, cardiac, line-noise, and channel artifacts that must be removed before any meaningful analysis~\cite{MutanenEtAl2018, BlumEtAl2019}. Over the past decade the community has responded by building automated preprocessing pipelines that scale to large datasets where manual cleaning is infeasible~\cite{BigdelyShamloEtAl2015, GabardDurnamEtAl2018, PedroniEtAl2019}. Yet a foundational problem has gone unsolved: there is no objective, task-independent way to decide whether one preprocessing run is better than another. As \citet{PedroniEtAl2019} state plainly, ``there is no method to objectively quantify the quality of preprocessed EEG,'' a void that pushes researchers toward subjective subject exclusion and inflates the risk of analytic flexibility. Compounding the problem, \citet{RobbinsEtAl2020} show that downstream EEG results shift substantially and non-uniformly with preprocessing choices, so the choice demonstrably matters even though it is currently made ad hoc.

The stakes of this gap are practical. Because preprocessing quality is judged either by downstream task accuracy --- which requires labels and a classifier --- or by subjective visual inspection, large public resting-state corpora cannot be quality-controlled in a principled, reproducible way. A label-free, recording-level quality metric would allow preprocessing to be \emph{optimized} rather than merely \emph{applied}, and would make the choice of pipeline parameters an empirical question rather than a matter of convention.

\subsection{Background and key concepts}

Two recent developments make a principled solution conceivable. The first is the parameterization of neural power spectra into separable \emph{periodic} and \emph{aperiodic} components~\cite{HallerEtAl2018, DonoghueEtAl2020b}. In this framework, a power spectrum is modeled as a $1/f$-like aperiodic background (an offset and an exponent) plus a small number of Gaussian-shaped oscillatory peaks. Crucially, the two components are physiologically distinct and can be conflated by naive band measures: \citet{DonoghueEtAl2020a} show that the canonical theta/beta band ratio is driven largely by aperiodic differences rather than oscillatory power. An alternative decomposition, Irregular-Resampling Auto-Spectral Analysis (IRASA), separates fractal from oscillatory power by resampling~\cite{WenLiu2016}, and \citet{GersterEtAl2022} provide a careful comparison of the two, quantifying exactly when each fails. This separability is the conceptual basis for treating oscillatory power as \emph{signal} and aperiodic power as a \emph{noise-like} background --- and hence for a spectral signal-to-noise ratio.

The second development is the maturation of automated pipeline optimization. Classical Automated Machine Learning (AutoML) frames algorithm and hyperparameter selection as a single search problem~\cite{ThorntonEtAl2012, FeurerEtAl2015, BischlEtAl2023}, and recent work extends this to data-side decisions such as cleaning~\cite{NeutatzEtAl2022, SiddiqiEtAl2023}. More recently still, large language model (LLM) agents that reason over intermediate state, call tools, and search have begun to match or surpass hand-tuned machine-learning engineering~\cite{GuoEtAl2024, HuangEtAl2023, ChiEtAl2024}. Together, these strands suggest that preprocessing could be cast as a search problem whose objective is a spectral quality metric and whose controller is an agent --- if the metric existed and the agent could be transferred to electrophysiology.

\subsection{State of the art}

\textbf{Automated preprocessing pipelines.} The field has converged on a canonical sequence of operations: robust referencing and bad-channel detection~\cite{BigdelyShamloEtAl2015}, filtering and line-noise removal, independent component analysis (ICA) with automated component classification~\cite{PionTonachiniEtAl2019, LiEtAl2022}, automated bad-trial rejection~\cite{JasEtAl2017}, and Artifact Subspace Reconstruction for transient artifacts~\cite{BlumEtAl2019}. These primitives are bundled into full pipelines for developmental, low-density, ERP, neonatal, and EEG-fMRI data~\cite{GabardDurnamEtAl2018, LopezEtAl2022, MonachinoEtAl2022, KumaravelEtAl2022, MayeliEtAl2021}, and into resting-state feature-extraction tools~\cite{vilaEtAl2023}. Each pipeline exposes a handful of tunable parameters, but applies them with fixed defaults rather than searching for a per-recording optimum.

\textbf{Spectral parameterization as a validated decomposition.} Beyond its descriptive use, the aperiodic component has been shown to carry genuine, behaviorally and clinically meaningful information: the $1/f$ slope tracks arousal and distinguishes wakefulness from anesthesia and sleep~\cite{LendnerEtAl2020}, and aperiodic parameters change systematically across development and aging~\cite{HillEtAl2022, MerkinEtAl2022} and dissociate from periodic activity in their links to cognition~\cite{ThuwalEtAl2021, WilsonEtAl2022, OstlundEtAl2022}. This validity is a double-edged sword for any quality metric: the aperiodic component cannot simply be minimized as ``noise'' without risking the destruction of real signal.

\textbf{Objective preprocessing quality and pipeline search.} A small but growing body of work asks how to judge preprocessing objectively. \citet{Delorme2023} builds an explicit quality metric --- the percentage of significant channels in a post-stimulus window --- and, applying it across optimized pipelines in four major toolboxes, finds the provocative result that most automated correction steps either do not help or actively hurt, with only high-pass filtering and bad-channel interpolation reliably beneficial. \citet{CoelliEtAl2023} manually compares methods at each preprocessing step and proposes quantitative indices for selection. On the optimization side, \citet{BashashatiEtAl2016} applied Bayesian optimization to per-user brain--computer-interface pipelines, and \citet{MartinezEtAl2023} search a per-instance preprocessing pipeline --- the closest prior art to adaptive, recording-specific preprocessing. None of this work, however, uses a spectral signal-to-noise objective, and none applies a modern agentic controller to EEG.

\subsection{Research gaps}

The field's achievements are real: preprocessing is now automatable, the periodic/aperiodic split is physiologically validated, and pipeline search is a mature methodology. Yet three coupled gaps remain, and they sit precisely at the unoccupied intersection of these strands. First, \emph{no study uses a FOOOF-derived signal-to-noise ratio as the objective function for preprocessing}; the literature evaluates preprocessing by task accuracy or by label-dependent contrast metrics~\cite{RobbinsEtAl2020, Delorme2023}, leaving the label-free, resting-state case open even though the enabling decomposition exists~\cite{DonoghueEtAl2020b}. Second, \emph{no exhaustive-search ground-truth resource exists for EEG preprocessing}: \citet{CoelliEtAl2023} perform the per-step comparison manually on a few methods, and \citet{MartinezEtAl2023} search pipelines but for non-EEG data optimizing task loss, so a per-recording optimal-configuration dataset under a spectral metric does not exist. Third, \emph{LLM agents have never been applied to neurophysiological preprocessing}: the state-of-the-art agentic ML-engineering systems are all evaluated on tabular, code, or chemistry tasks~\cite{GuoEtAl2024, ChiEtAl2024, TriratEtAl2024, JiangEtAl2025}, never on raw electrophysiology. A fourth consideration --- that a naive spectral SNR could penalize physiologically real aperiodic activity~\cite{LendnerEtAl2020, GersterEtAl2022} --- is not a separate gap but a validity constraint that any such metric must internalize.

\subsection{Hypothesis and approach}

We propose and validate an agentic EEG preprocessing pipeline that unifies these three gaps into one program. The central hypothesis is that \emph{optimal, per-recording preprocessing configurations selected by maximizing a FOOOF-derived signal-to-noise ratio yield significantly higher spectral quality than fixed standard pipelines, and that an LLM agent can recover near-optimal configurations from intermediate signal features}. Concretely, we (i) define a FOOOF-SNR that treats periodic power as signal and aperiodic-plus-residual power as noise, guarded against the aperiodic-as-signal pitfall by band-limited scoring and goodness-of-fit gating; (ii) build a ground-truth oracle by exhaustively searching a preprocessing parameter grid over public datasets and recording the best configuration per recording; and (iii) train and evaluate an LLM-agent controller that selects preprocessing operations and parameters, using the oracle as both supervision and an upper bound.

\subsection{Contributions}

This work makes the following contributions:
\begin{itemize}
  \item A \textbf{label-free, task-independent quality metric} (FOOOF-SNR) for EEG preprocessing, validated against exhaustive search and expert ratings, with an explicit safeguard against penalizing physiologically meaningful aperiodic activity.
  \item A \textbf{reusable ground-truth benchmark} of per-recording optimal preprocessing configurations, constructed by exhaustive/greedy search over public EEG datasets under the FOOOF-SNR objective.
  \item An \textbf{LLM-agent preprocessing controller} --- the first application of agentic, tool-using LLMs to electrophysiological preprocessing --- evaluated against fixed pipelines and classical search baselines for quality, cost, and generalization.
  \item An \textbf{interpretability analysis} of the agent's learned policy, testing whether it recovers established preprocessing heuristics and surfaces recording-specific rules.
\end{itemize}

The remainder of this paper is organized as follows. Section~\ref{sec:related} reviews related work across automated preprocessing, spectral parameterization, and pipeline optimization. Section~\ref{sec:method} describes the FOOOF-SNR metric, the oracle construction, and the agent. Section~\ref{sec:results} reports the experiments, and Section~\ref{sec:discussion} discusses implications and limitations.
